% Part: first-order-logic
% Chapter: completeness
% Section: completeness-theorem

\documentclass[../../../include/open-logic-section]{subfiles}

\begin{document}

\iftag{FOL}
      {\olfileid[zh]{fol}{com}{cth}}
      {\olfileid[zh]{pl}{com}{cth}}

\olsection{完全性定理}

\begin{explain}
让我们把我们的结果结合起来：我们得到了完全性定理。
\end{explain}

\begin{thm}[完全性定理]
\ollabel{thm:completeness}
设 $\Gamma$ 是一组!!{sentence}。如果 $\Gamma$ 是一致的，它就是可满足的。
\end{thm}

\begin{proof}
设 $\Gamma$ 是一致的。\iftag{FOL}{由
  \olref[hen]{lem:henkin}，存在一个饱和的、一致的集合
  $\Gamma' \supseteq \Gamma$。}{} 由 \olref[lin]{lem:lindenbaum}，存在
一个 $\Gamma^* \supseteq \iftag{FOL}{\Gamma'}{\Gamma}$，它是一致的且!!{complete}。
\iftag{FOL}{
  由于 $\Gamma' \subseteq \Gamma^*$，对每个!!{formula}~$!A(x)$，$\Gamma^*$ 都含有如下形式的!!a{sentence}：
  \iftag{prvEx}
        {$\lexists[x][!A(x)] \lif !A(c)$}
        {$\lnot\lforall[x][!A(x)] \lif \lnot !A(c)$}
  因而 $\Gamma^*$ 是饱和的。如果 $\Gamma$
  不含有~$\eq$，那么}{}
由\olref[mod]{lem:truth}，
$\iftag{FOL}{\Sat{M(\Gamma^*)}}{\pSat{v(\Gamma^*)}}{!A}$ 当且仅当 $!A \in
\Gamma^*$。由此特别可得，对所有 $!A \in
\Gamma$，$\iftag{FOL}{\Sat{M(\Gamma^*)}}{\pSat{v(\Gamma^*)}}{!A}$，所以
$\Gamma$ 是可满足的。\iftag{FOL}{
  如果 $\Gamma$ 含有~$\eq$，
  那么由 \olref[ide]{lem:truth}，对所有!!{sentence}~$!A$，
  $\Sat{\equivclass{M}{\approx}}{!A}$ 当且仅当 $!A \in \Gamma^*$。特别地，
  $\Sat{\equivclass{M}{\approx}}{!A}$ 对所有 $!A \in
  \Gamma$ 成立，所以 $\Gamma$ 是可满足的。}{}
\end{proof}

\begin{cor}[完全性定理，第二形式]
\ollabel{cor:completeness}
对所有 $\Gamma$ 和!!{sentence}~$!A$：如果 $\Gamma \Entails !A$，则
$\Gamma \Proves !A$。
\end{cor}

\begin{proof}
注意 \olref{cor:completeness} 和 \olref{thm:completeness} 中的
$\Gamma$ 是全称量化的。为了不把我们自己弄糊涂，让我们换用不同的变元来重新表述 \olref{thm:completeness}：对任何一组!!{sentence}~$\Delta$，如果
$\Delta$ 是一致的，它就是可满足的。由逆否，如果
$\Delta$ 是不可满足的，那么 $\Delta$ 就是不一致的。我们将用这个来证明该推论。

设 $\Gamma \Entails !A$。那么根据 \olref[syn][sem]{prop:entails-unsat}，$\Gamma \cup \{\lnot !A\}$ 是
不可满足的。把 $\Gamma
\cup \{\lnot !A\}$ 作为我们的~$\Delta$，前面的
\olref{thm:completeness} 版本给出 $\Gamma \cup \{\lnot !A\}$ 是
不一致的。根据
\iftag{FOL}{%
  \tagrefs{prfAX/{fol:axd:prv:prop:prov-incons},
    prfSC/{fol:seq:prv:prop:prov-incons},
    prfND/{fol:ntd:prv:prop:prov-incons},
    prfTab/{fol:tab:prv:prop:prov-incons}}}{%
  \tagrefs{prfAX/{pl:axd:prv:prop:prov-incons},
    prfSC/{pl:seq:prv:prop:prov-incons},
    prfND/{pl:ntd:prv:prop:prov-incons},
    prfTab/{pl:tab:prv:prop:prov-incons}}},
$\Gamma \Proves !A$。
\end{proof}

\tagprob{FOL}
\begin{prob}
用 \olref[fol][com][cth]{cor:completeness} 证明
\olref[fol][com][cth]{thm:completeness}，从而表明完全性定理的两种
表述是等价的。
\end{prob}
\tagendprob

\tagprob{notFOL}
\begin{prob}
用 \olref[pl][com][cth]{cor:completeness} 证明
\olref[pl][com][cth]{thm:completeness}，从而表明完全性定理的两种
表述是等价的。
\end{prob}
\tagendprob

\tagprob{FOL}
\begin{prob}
为了使!!a{derivation}系统是完全的，它的规则必须
强大到足以把所有不可满足的集合证明为不一致。在证明完全性的过程中，哪些!!{derivation}规则是必需的？这些规则中有没有
哪条在证明中根本没用上？为了回答
这些问题，请列出清单或画出图表，说明哪些!!{derivation}规则被用于哪些通向
\olref[fol][com][cth]{thm:completeness} 的证明的结果中。请务必记下这些证明中对规则的任何默示
使用。
\end{prob}
\tagendprob

\tagprob{notFOL}
\begin{prob}
为了使!!a{derivation}系统是完全的，它的规则必须
强大到足以把所有不可满足的集合证明为不一致。在证明完全性的过程中，哪些!!{derivation}规则是必需的？这些规则中有没有
哪条在证明中根本没用上？为了回答
这些问题，请列出清单或画出图表，说明哪些!!{derivation}规则被用于哪些通向
\olref[pl][com][cth]{thm:completeness} 的证明的结果中。请务必记下这些证明中对规则的任何默示
使用。
\end{prob}
\tagendprob

\end{document}
